\section{Discussion and limitations}
\label{sec:discussion}

\subsection{Where the Sudoku logic shows up}
\label{sec:discussion-summary}

The three confirmatory results, summarized in Section~\ref{sec:summary}, give a consistent, if partly negative, picture of one coarse-to-fine Sudoku generator: under the frequency policy one replaced level-1 code cost far more than one replaced level-11 code, while levels 2 and 11 were practically indistinguishable; regenerating the later levels attenuated that damage by an amount whose practical relevance is not established; and the candidate mask was linearly readable from the level-6 state only weakly, below the registered practical threshold, while the level-12 readout was no better than from the teacher state that encodes the solution and, in an exploratory analysis, not equivariant under the sampled symmetries.

Taken together, the results are consistent with a hierarchy in which a coarse level carries something the solution depends on, while the final decoder recovers most of the cell content only from the late levels (Section~\ref{sec:e1}); part-way through, the candidate structure of the puzzle is linearly readable only weakly, below the practical threshold. Each part of this picture is conditional on the checkpoint, the replacement policy, the probe family and its training budget. In particular, none of the results shows that the coarse code is a plan of a solving algorithm, that the model repairs logical errors, or that the intermediate state lacks candidate information that another readout could find.

\subsection{Relation to prior work}
\label{sec:discussion-related}

Linear readouts of Sudoku candidates, and a causally verified circuit for naked singles, have been reported for an autoregressive transformer trained on solver traces \citep{kniazev2026linear}; in that work both per-cell candidate probes and unit-level digit-placement probes reached perfect exact match in the middle layers, while per-cell probes of the filled digits did not, and the authors conclude that the board state is organized by rows, columns and boxes rather than by cells. In a hierarchical reasoning model, probe directions for Sudoku were decodable but not more causal than random directions \citep{rodrigues2026hrm}. Our setting differs in the object of study: a non-autoregressive generator whose intermediate states are built from discrete, level-indexed codes. In this setting the candidate mask was only weakly linearly readable at level~6, and our pooled readout of the unit-level target was at chance on the generated level-6 state, while both were readable from the input. Whether the difference comes from the architecture, the training objective or the readout cannot be decided from one checkpoint. Exploratory notes on the released NCM Sudoku checkpoints, written with substantial AI assistance and not peer reviewed, report a naked-single probe on final-level cell states that is decodable above a shuffled-label control but whose causal ablation was not specific, and residual-omission and state-exchange tests under fixed and regenerated continuations whose criteria were not met \citep{quantix2026ncmnotes}; our study differs by its preregistration, the equal-cost contrast of single-code replacements at levels 1 and 11, the candidate-mask target on the intermediate cumulative state and its control matrix. Generated NCM codes for text have been reported, in a preprint that is not peer reviewed, to be largely non-canonical under re-encoding, while for the Sudoku checkpoint the same preprint reports far fewer rewritten codes per level \citep{khoury2026roundtrip}; because different code sequences can decode to the same grid, we report cell-level outcomes next to every code-level count. The compensation we measure is related to, and kept distinct from, self-repair within a single forward pass \citep{ahmad2026dose}. In masked diffusion language models, remasking an isolated error on Sudoku has been reported to usually reproduce it when the surrounding context stays fixed \citep{oh2026regenerate}, a behavioral observation, in a different architecture, that regenerating part of a solution under a fixed erroneous context need not undo the error. Candidate-set evidence in other work on logic puzzles comes from forcing outputs and comparing top-ranked predictions with a solver rather than from a learned probe \citep{shah2024logic}. The probing methodology follows the control-task tradition \citep{hewitt2019control,belinkov2022probing}.

\subsection{Limitations}
\label{sec:discussion-limitations}

\paragraph{One checkpoint, one task, one policy family.} All results concern the released Sudoku-hard checkpoint and our own puzzle distribution, and the E2 results are conditional on the code replacement policies. The training data of the checkpoint is not released, so overlap with our puzzles cannot be excluded; the distribution check only shows that our puzzles are about as hard for the model as the released evaluation set.

\paragraph{Distribution shift after replacement.} A replaced code need not be natural given the puzzle and the earlier levels, and the H2 contrast measures an equal number of discrete edits, not an equal semantic change; level~1 is architecturally distinct from the other levels. Replacements can produce format errors, \R{e2.l1.uniform.fmtC}\% of the level-1 replacements in branch C under the uniform policy against \R{e2.l1.frequency.fmtC}\% under the frequency policy, and such outputs stay in the solve denominator as unsolved (Appendix~\ref{sec:appendix-limits}).

\paragraph{What the cumulative state is built for.} The cumulative state is the sum of codebook vectors of a tokenizer trained to compress solved grids, and the generator receives the puzzle at every level. Candidate reasoning, if the generator performs any, could therefore happen inside the generator's activations before it emits a level's codes and leave little trace in the cumulative state; the weak level-6 readout is consistent with this, but our design does not test it.

\paragraph{Trained robustness.} The NCM paper describes training the generator with corruption of earlier levels \citep{kumbong2026ncm}; assuming the released Sudoku checkpoint was trained the same way, the compensation in H3 is consistent with what training rewards and is not evidence of spontaneous error handling.

\paragraph{Readout capacity.} The primary E3 probes are linear, trained with a fixed budget, and the training-size curve had not leveled off; a nonlinear readout of the cell together with its peers, peer information beyond the fixed principal-component projection, larger training sets and other representations were not tested (a per-cell MLP improved only marginally), and $\theta_1$ is a gain over the permuted control, not over a trivial predictor (Appendix~\ref{sec:appendix-limits}).

\paragraph{Teacher and generated states.} At level~12 the generated and the teacher state were read equally well, so the level-12 readability cannot be separated from solution content. The same holds at level~6: the teacher's level-6 state, computed from the solution rather than the puzzle, reached the same Brier loss as the generated state (\R{brier.teacher.L06} against \R{brier.wl.gen.L06}), so the small gain of the generated state over the permuted control is also present in a state that carries no clue information and need not reflect the puzzle's candidate structure. At level~12, a solution-digit-only reference, computed after TEST from the TEST candidate rate, gave a Brier loss of about \R{soldigit.brier}, close to both level-12 probes (exploratory; Appendix~\ref{sec:appendix-limits}).

\paragraph{Design changes and post-registration decisions.} We changed three design choices before the lock after pilot results on non-TEST data (Section~\ref{sec:setup-stats}); we moved the H2 and H3 levels to where the pilot had shown an effect. After the lock, no hypothesis, threshold or sample size changed; the remaining decisions are listed in Appendix~\ref{sec:appendix-deviations}, and the registered dose-curve comparison could not be computed because the TEST damage values did not cover the frozen range. Two analyses that the registration lists as exploratory were not run and are not reported: the supplementary generalized linear mixed model, whose formula the protocol had left open, and the teacher control with different clue masks of the same solution, for which no recipe had been fixed.

\paragraph{The causal study was not run.} The locked protocol described a conditional targeted patching study (E4), to be run only if a reproducible, interpretable logical readout and acceptable equivariance were found on a separate selection set; the registration excluded E4 and left it for a separate registration before its test set was generated. Given the level-6 readout below the practical threshold and the non-equivariant level-12 readout, we judged the entry condition very unlikely to hold. We dropped the study without generating its selection or test sets, and we made this decision after seeing the TEST results of E3; it does not change any confirmatory result, but it means that we make no causal claim about the readout directions here.
